\documentclass[letterpaper,11pt]{article}
\usepackage[T1]{fontenc}
\usepackage[utf8]{inputenc}
\usepackage[margin=1in]{geometry}
\usepackage[hidelinks]{hyperref}
\usepackage{xcolor}
\usepackage[default]{sourcesanspro}
\usepackage{parskip}
\usepackage[protrusion=true,expansion=false]{microtype}
%----------COLORS----------
\definecolor{ResumeNavy}{HTML}{18324A}
\definecolor{ResumeSlate}{HTML}{5B6573}
\definecolor{ResumeText}{HTML}{111827}
\color{ResumeText}
\setlength{\parskip}{8pt}
\pagestyle{empty}
\begin{document}
%----------HEADER (matches resume identity)----------
\begin{center}
  {\LARGE \bfseries \textcolor{ResumeNavy}{Deva Anand}} \\ \vspace{3pt}
  \small
  (413) 315-1715 \enspace\textbar\enspace
  \href{mailto:devaanand@umass.edu}{devaanand@umass.edu} \enspace\textbar\enspace
  \href{https://devaanand.com}{devaanand.com} \enspace\textbar\enspace
  \href{https://linkedin.com/in/devaaa/}{linkedin.com/in/devaaa} \enspace\textbar\enspace
  \href{https://github.com/Deva-1903}{github.com/Deva-1903}
\end{center}
\vspace{4pt}
{\color{ResumeNavy}\hrule height 1pt}
\vspace{12pt}

May 24, 2026

Snowflake University Recruiting \\
{\color{ResumeSlate}\textit{Re: Software Engineer Intern (Infrastructure Automation) --- Fall 2026}}

\vspace{4pt}

Dear Snowflake Hiring Team,

I am a Master's student in Computer Science at UMass Amherst (graduating May 2027), applying for your Fall 2026 Software Engineer Internship in Infrastructure Automation. What draws me to this role is its leverage: building the automation frameworks, tooling, and release infrastructure that let everyone else ship faster and more reliably. I would rather make a whole engineering org more productive than ship one more feature, and a data platform that real customers depend on is exactly where that work matters most.

That instinct shows up in what I have already built. At Wysa I built an internal annotation platform that automated a manual, spreadsheet-based workflow into secure role-based APIs, cutting 100+ hours of repetitive operations a month and lifting throughput 60\% for the team that relied on it. I also owned reliability for a multi-tenant NHS backend serving 8+ UK clients, integration tests, authentication hardening, 20+ VAPT remediations, and release monitoring, because trustworthy software is the whole point on sensitive data. More recently I optimized a Spark ETL pipeline over a 1.4-billion-row dataset: profiling stage-level timings in the Spark UI, enabling adaptive query execution and skew-join handling, cutting runtime by roughly 25\%, and turning the fixes into a repeatable optimization playbook rather than a one-off. That mix of large-scale data, performance, and repeatable automation maps closely to your Engineering Systems and Quality \& Release Engineering work.

Your description of treating AI as a high-trust collaborator matches how I work day to day. I lean on tools like Claude Code and Cursor as accelerators, but I keep system behavior grounded in tests, CI, and verifiable evidence. My agentic web-discovery project, for instance, pairs LLM extraction with deterministic validation, cell-level provenance, and 150+ automated tests. I have not used Jenkins, Ansible, or Terraform in production yet, but I learn fastest when I am a little out of my depth, and picking up new infrastructure tooling quickly is exactly the kind of challenge I am looking for. I would be glad to bring that to a real system Snowflake customers rely on.

Thank you for your time and consideration.

\vspace{6pt}
Sincerely, \\
\textbf{Deva Anand}

\end{document}
